\documentclass[12pt]{article}

\usepackage{url}
\usepackage{geometry}
\geometry{a4paper}
\usepackage{microtype}

\usepackage{graphicx}
\usepackage{amsmath}
\usepackage{algorithm}
\usepackage[noend]{algpseudocode}
\usepackage{float}
\usepackage{wrapfig}
\usepackage{listings}
\usepackage{multirow}

\usepackage{enumitem}

%para escribir en español
\usepackage[utf8x]{inputenc}
\usepackage[spanish]{babel}

\renewcommand{\contentsname}{Índice}
\renewcommand{\refname}{Referencias}
\renewcommand{\figurename}{Figura}

\usepackage[nottoc,numbib]{tocbibind}

\linespread{1.2}

\graphicspath{{./figures/}} % Specifies the directory where pictures are stored

\begin{document}

\begin{titlepage}

    \center

    \textsc{\LARGE Universidad de la República}\\[1.5cm]
    \textsc{\Large Facultad de Ingeniería}\\[1.0cm]
    \textsc{\large Recuperación de Información y Recomendaciones en la Web}\\[0.5cm]

    \rule{\linewidth}{0.5mm} \\[1cm]
    {\huge \bfseries Amazon Copilot}\\[0.5cm]
    \rule{\linewidth}{0.5mm} \\[1.5cm]

    \begin{minipage}{0.4\textwidth}
        \begin{flushleft} \large
            \emph{Autores:}\\
            Juan Pablo Conde\\
            Xavier Iribarnegaray\\
            Juan Pablo Sotelo\\
        \end{flushleft}
    \end{minipage}
    ~
    \begin{minipage}{0.4\textwidth}
        \begin{flushright} \large
            \emph{Profesores:} \\
            Libertad Tansini \\
        \end{flushright}
    \end{minipage}\\[2cm]

    {\large \today}\\[1.5cm]

    \includegraphics[width=0.4\textwidth]{fing-logo}\\[1cm]

    \vfill

\end{titlepage}


\pagenumbering{arabic}

\tableofcontents

\newpage

\section{Introducción}
En este documento presentamos la propuesta para Amazon Copilot, un proyecto que buscará facilitar la búsqueda y selección de productos en plataformas de e-commerce. Nuestro sistema planea utilizar un dataset de productos de Amazon para implementar tres funcionalidades clave:

\begin{itemize}
    \item \textbf{Búsqueda híbrida de productos:} Combinaremos búsqueda semántica y técnicas de text matching tradicionales para ofrecer resultados más relevantes a las consultas de los usuarios.

    \item \textbf{AI Agent:} Un asistente conversacional que utilizará las capacidades
          de búsqueda implementadas y podrá refinar los requerimientos del usuario mediante
          diálogo natural.

    \item \textbf{Sistema de recomendación:} Ofrecerá sugerencias de productos relacionados con los ítems seleccionados en el carrito de compras.
\end{itemize}

El proyecto integrará técnicas modernas de procesamiento de lenguaje natural, recuperación de información y sistemas de recomendación para crear una experiencia de usuario fluida y personalizada.

\newpage
\section{Arquitectura del Sistema}

\section{Funcionalidades}

\section{Implementación Técnica}

\section{Referencias}

\end{document}
